\chapter{Open Quantum Systems and Nonlinear PDEs}
\label{ch:lec14}

\begin{objectives}
\begin{itemize}[leftmargin=1.4em,nosep]
  \item Describe Lindblad dynamics and its dilation to a unitary.
  \item Prepare thermal and ground states dissipatively.
  \item Outline how nonlinear PDEs are linearized (Carleman, Cole--Hopf) for quantum solvers.
\end{itemize}
\end{objectives}

The last three chapters kept the dynamics \emph{linear}: the Schr\"odinger equation
(\cref{ch:lec11,ch:lec12}) and the general linear ODE $\dot x=Ax+b$ (\cref{ch:lec13}). This
chapter relaxes both idealizations, in two independent directions that nonetheless share a common
strategy---\emph{embed the hard problem inside a larger linear or unitary one}. The first half
treats \emph{open} quantum systems, whose state is a density operator $\rho$ evolving under a
non-unitary but completely positive semigroup; the governing equation is the Lindblad master
equation, the continuous-time analogue of the CPTP channels of \cref{ch:lec02}. We simulate it
by dilating its non-unitary generator to a unitary on a larger space (Stinespring,
\cref{thm:lec02-stinespring}) and by operator splitting and a truncated Dyson series, all
reusing the Hamiltonian simulation of \cref{ch:lec11,ch:lec12}. We then turn dissipation into a
resource: by engineering the jump operators so that a target thermal or ground state is the
\emph{unique fixed point}, Lindblad dynamics prepares that state, an open-system cousin of the
adiabatic and Zeno solvers of \cref{ch:lec10}. The second half attacks \emph{nonlinear} PDEs.
Because quantum mechanics is linear, a nonlinear evolution must first be \emph{linearized}: we
present Carleman linearization, which embeds a quadratic ODE into an infinite linear system on
the tensor powers of the state and truncates it (solving the result with the methods of
\cref{ch:lec13}), and the exact Cole--Hopf transformation, which maps the viscous Burgers
equation onto the heat equation. Throughout, $\norm\cdot$ is the spectral norm, $\norm\cdot_1$
the trace norm, $\norm\cdot_\diamond$ the diamond norm of \cref{ch:lec02}, and $\mathcal D(\Hilb)$
the set of density operators on $\Hilb$.

\section{Lindblad dynamics}
\label{sec:lec14-1}

\paragraph{The master equation.}
An open system interacts with environmental degrees of freedom that are not modeled explicitly.
In the Markovian regime---environmental correlations decay fast on the system's timescale---the
state $\rho(t)\in\mathcal D(\C^d)$ obeys the Gorini--Kossakowski--Sudarshan--Lindblad (GKSL)
master equation
\begin{equation}
  \frac{\dd\rho}{\dd t}
  =\underbrace{-\ii[H,\rho]}_{\mathcal L_H(\rho)}
   +\underbrace{\sum_{j=1}^{J}\Bigl(L_j\rho L_j\dg-\tfrac12\bigl\{L_j\dg L_j,\rho\bigr\}\Bigr)}
     _{\mathcal L_D(\rho)}
  =:\mathcal L(\rho),
  \label{eq:lec14-lindblad}
\end{equation}
where $H=H\dg$ is the system Hamiltonian, the $L_j\in\C^{d\times d}$ are \emph{jump operators}
encoding the coupling to the environment, $[A,B]=AB-BA$, and $\{A,B\}=AB+BA$. The generator
$\mathcal L$ is the \emph{Lindbladian}; its coherent part $\mathcal L_H$ reproduces the
Liouville--von Neumann equation, while its dissipative part $\mathcal L_D$ is genuinely
irreversible. The formal solution is the one-parameter semigroup
\begin{equation}
  \rho(t)=\Phi_t\bigl(\rho(0)\bigr),\qquad \Phi_t:=\ee^{\mathcal L t},\quad t\ge0.
  \label{eq:lec14-semigroup}
\end{equation}
When all $L_j=0$ this collapses to $\rho(t)=\ee^{-\ii Ht}\rho(0)\ee^{\ii Ht}$, closed-system
unitary evolution. The value of the Lindblad form is that $\Phi_t$ is always a legitimate
channel.

\begin{theorem}[$\ee^{\mathcal L t}$ is a quantum channel]
\label{thm:lec14-cptp}
For every $t\ge0$ the map $\Phi_t=\ee^{\mathcal L t}:L(\C^d)\to L(\C^d)$ defined by
\eqref{eq:lec14-lindblad}--\eqref{eq:lec14-semigroup} is completely positive and trace preserving
(CPTP), i.e.\ a quantum channel in the sense of \cref{def:lec02-channel}. The family
$\{\Phi_t\}_{t\ge0}$ forms a semigroup: $\Phi_0=\mathrm{id}$ and $\Phi_{t+s}=\Phi_t\Phi_s$.
\end{theorem}

\begin{proof}
The semigroup laws are immediate from $\ee^{\mathcal L\cdot0}=\mathrm{id}$ and
$\ee^{\mathcal L(t+s)}=\ee^{\mathcal L t}\ee^{\mathcal L s}$. For the CPTP claim, fix a small step
$\Delta t=t/N_t$ and construct a genuine channel $\Phi$ with
$\norm{\ee^{\mathcal L\Delta t}-\Phi}_\diamond\le C_1\Delta t^2$; we exhibit such a $\Phi$
explicitly in \cref{sec:lec14-2} (equation \eqref{eq:lec14-dilated-channel}). Granting this, let
$C_2=\sup_{0\le s\le t}\norm{\ee^{\mathcal L s}}_\diamond$. Because the diamond norm is
submultiplicative and $\norm\Phi_\diamond=1$, the telescoping identity
\[
  \ee^{\mathcal L t}-\Phi^{N_t}
  =\sum_{k=0}^{N_t-1}\ee^{\mathcal L(N_t-1-k)\Delta t}\,
     \bigl(\ee^{\mathcal L\Delta t}-\Phi\bigr)\,\Phi^{k}
\]
gives $\norm{\ee^{\mathcal L t}-\Phi^{N_t}}_\diamond\le C_2 C_1 N_t\Delta t^2=C_1C_2\,t\,\Delta t\to0$
as $N_t\to\infty$. Thus $\ee^{\mathcal L t}$ is a diamond-norm limit of the channels
$\Phi^{N_t}$ (a composition of channels is a channel). Since the set of CPTP maps is closed in the
diamond norm, $\ee^{\mathcal L t}$ is itself CPTP.
\end{proof}

The converse is the structural GKSL theorem: in finite dimension, \eqref{eq:lec14-lindblad} is
the \emph{most general} generator of a norm-continuous semigroup of channels.

\begin{proposition}[GKSL characterization]
\label{prop:lec14-gksl}
Let $\{\Phi_t\}_{t\ge0}$ be a norm-continuous semigroup of quantum channels on $L(\C^d)$ with
generator $\mathcal G(\rho)=\lim_{t\to0^+}(\Phi_t(\rho)-\rho)/t$. Then there exist a Hermitian
$H$ and operators $L_1,\dots,L_J$ such that $\mathcal G$ has the form of the right-hand side of
\eqref{eq:lec14-lindblad}.
\end{proposition}

\begin{proof}[Proof sketch]
For small $t$ the channel $\Phi_t$ has a Kraus form $\Phi_t(\rho)=\sum_\alpha K_\alpha(t)\rho
K_\alpha(t)\dg$ with $\sum_\alpha K_\alpha(t)\dg K_\alpha(t)=\id$. Expanding the Choi matrix to
first order in $t$, one may choose the Kraus operators so that one of them is close to the
identity, $K_0(t)=\id+t\mathcal G_0+o(t)$, while the rest are $\mathcal O(\sqrt t)$, namely
$K_j(t)=\sqrt t\,L_j+\mathcal O(t)$. The trace-preservation condition, expanded to first order,
forces $\mathcal G_0+\mathcal G_0\dg+\sum_{j}L_j\dg L_j=0$, so $\mathcal G_0=-\ii H-\tfrac12
\sum_j L_j\dg L_j$ for a Hermitian $H$. Substituting back and collecting the $\mathcal O(t)$
terms of $(\Phi_t(\rho)-\rho)/t$ reproduces \eqref{eq:lec14-lindblad}.
\end{proof}

\paragraph{The adjoint (Heisenberg) equation.}
Expectation values of an observable $O$ evolve through the \emph{adjoint} Lindbladian
$\mathcal L\dg$, defined by $\Tr[\mathcal L(\rho)\,O]=\Tr[\rho\,\mathcal L\dg(O)]$ under the
Hilbert--Schmidt pairing $\braket{X|Y}=\Tr[X\dg Y]$:
\begin{equation}
  \frac{\dd}{\dd t}\Tr[\rho(t)O]=\Tr[\rho(t)\,\mathcal L\dg(O)],\qquad
  \mathcal L\dg(O)=\ii[H,O]+\sum_{j=1}^{J}\Bigl(L_j\dg O L_j-\tfrac12\{L_j\dg L_j,O\}\Bigr).
  \label{eq:lec14-adjoint}
\end{equation}
The Heisenberg picture $O(t)=\ee^{\mathcal L\dg t}(O)$ then satisfies $\Tr[\rho(t)O]=
\Tr[\rho(0)O(t)]$, exactly as for closed systems.

\begin{example}[{Classical Markov chains as Lindblad dynamics}]
\label{exa:lec14-markov}
Take $H=0$ and, for a rate matrix $q_{ij}\ge0$ ($i\ne j$), jump operators
$L_{ij}=\sqrt{q_{ij}}\,\ket{i}\!\bra{j}$. For a diagonal (classical) state
$\rho=\sum_j p_j\proj j$ one has $L_{ij}\rho L_{ij}\dg=q_{ij}p_j\proj i$ and $L_{ij}\dg L_{ij}
=q_{ij}\proj j$, so \eqref{eq:lec14-lindblad} keeps $\rho$ diagonal and the populations obey
$\dot p_i=\sum_{j\ne i}q_{ij}p_j-\bigl(\sum_{k\ne i}q_{ki}\bigr)p_i$, i.e.\ $\dot p=Qp$ with the
classical generator $Q_{ij}=q_{ij}$ ($i\ne j$), $Q_{jj}=-\sum_{i\ne j}q_{ij}$. Restricted to
diagonal states the Lindblad semigroup $\ee^{\mathcal L t}$ is exactly the column-stochastic
Markov semigroup $\ee^{Qt}$: Lindblad dynamics is the quantum analogue of a continuous-time
Markov chain.
\end{example}

\begin{example}[{Photon loss in a cavity}]
\label{exa:lec14-photon}
A single cavity mode with annihilation operator $a$, Hamiltonian $H=\omega a\dg a$, and a single
jump operator $L=\sqrt\kappa\,a$ models photon leakage at rate $\kappa$. With $n=a\dg a$ the
adjoint \eqref{eq:lec14-adjoint} gives $\mathcal L\dg(n)=\kappa(a\dg a\dg a a-n^2)=\kappa(n(n-1)
-n^2)=-\kappa n$, so $\dd\langle n\rangle/\dd t=-\kappa\langle n\rangle$ and $\langle n\rangle(t)
=\langle n\rangle(0)\ee^{-\kappa t}$: the mean photon number decays exponentially and the vacuum
$\proj0$ is the steady state.
\end{example}

\section{Dilation to a larger unitary system}
\label{sec:lec14-2}

\paragraph{The dilated Hamiltonian.}
The obstruction to simulating \eqref{eq:lec14-lindblad} directly is that $\Phi_t$ is not
unitary. Stinespring's theorem (\cref{thm:lec02-stinespring}) says every channel is a unitary on
a larger space followed by discarding an ancilla; we now realize the \emph{short-time}
Lindblad channel this way, using an ancilla of $a=\lceil\log_2(J+1)\rceil$ qubits with basis
$\{\ket0,\ket1,\dots,\ket J\}$. Define the \emph{dilated Hamiltonian}
\begin{equation}
  \widetilde H
  =\proj0\otimes\frac{H}{\sqrt{\Delta t}}
   +\sum_{j=1}^{J}\bigl(\ket j\!\bra0\otimes L_j+\ket0\!\bra j\otimes L_j\dg\bigr)
  =\begin{psmallmatrix}
     H/\sqrt{\Delta t} & L_1\dg & \cdots & L_J\dg\\
     L_1 & 0 & \cdots & 0\\
     \vdots & \vdots & & \vdots\\
     L_J & 0 & \cdots & 0
   \end{psmallmatrix},
  \label{eq:lec14-dilated}
\end{equation}
and set the short-time unitary $U=\ee^{-\ii\widetilde H\sqrt{\Delta t}}$. The corresponding
channel on the system, obtained by initializing the ancilla in $\ket0$ and tracing it out, is
\begin{equation}
  \Phi(\rho)=\Tr_a\!\Bigl(U\bigl(\proj0\otimes\rho\bigr)U\dg\Bigr).
  \label{eq:lec14-dilated-channel}
\end{equation}
The factors of $\sqrt{\Delta t}$ are chosen precisely so that the dissipative terms enter at the
right order. Writing $\widetilde H=\proj0\otimes H/\sqrt{\Delta t}+R$ with $R=\ket j\!\bra0
\otimes L_j+\mathrm{h.c.}$, one has $R^2=\mathrm{diag}(\sum_j L_j\dg L_j,\dots)$, so
$\norm{\widetilde H}^2=\mathcal O(\norm H^2\Delta t+\sum_j\norm{L_j}^2)$ stays bounded once
$\Delta t=\mathcal O(\norm{\mathcal L}_{\mathrm{be}}^{-1})$, where
\begin{equation}
  \norm{\mathcal L}_{\mathrm{be}}:=1+\norm H+\sum_{j=1}^{J}\norm{L_j}^2.
  \label{eq:lec14-Lbe}
\end{equation}

\begin{theorem}[Dilated-Hamiltonian simulation of Lindblad dynamics]
\label{thm:lec14-dilation}
Let $\Phi$ be the channel \eqref{eq:lec14-dilated-channel}, $\norm{\mathcal L}_{\mathrm{be}}$ as
in \eqref{eq:lec14-Lbe}, and $\Delta t=\mathcal O(\norm{\mathcal L}_{\mathrm{be}}^{-1})$. Then
the one-step error is $\norm{\ee^{\mathcal L\Delta t}-\Phi}_\diamond=\mathcal O
(\norm{\mathcal L}_{\mathrm{be}}^2\,\Delta t^2)$, and hence with $t=N_t\Delta t$,
\begin{equation}
  \norm{\ee^{\mathcal L t}-\Phi^{N_t}}_\diamond
  \le N_t\,\norm{\ee^{\mathcal L\Delta t}-\Phi}_\diamond
  =\mathcal O\!\bigl(t\,\norm{\mathcal L}_{\mathrm{be}}^2\,\Delta t\bigr).
  \label{eq:lec14-dilation-error}
\end{equation}
Consequently, achieving diamond-norm error $\epsilon$ at time $t$ requires
$N_t=\mathcal O(t^2\norm{\mathcal L}_{\mathrm{be}}^2/\epsilon)$ applications of the Hamiltonian
simulation $U=\ee^{-\ii\widetilde H\sqrt{\Delta t}}$, each followed by resetting the ancilla to
$\ket0$.
\end{theorem}

\begin{proof}[Proof sketch]
Expand $U=\ee^{-\ii\widetilde H\sqrt{\Delta t}}$ to third order in $\sqrt{\Delta t}$ as $W=\id
-\ii\widetilde H\sqrt{\Delta t}-\tfrac12\widetilde H^2\Delta t+\tfrac{\ii}6\widetilde H^3
\Delta t^{3/2}$, with $\norm{U-W}=\mathcal O(\norm{\widetilde H}^4\Delta t^2)=\mathcal O
(\norm{\mathcal L}_{\mathrm{be}}^2\Delta t^2)$ by the bound $\norm{\widetilde H}^2=\mathcal O
(\norm{\mathcal L}_{\mathrm{be}})$ above. Because $\norm\cdot_\diamond$ is defined over ancillary
extensions of trace-norm-one inputs and the partial trace is a contraction in the trace norm,
replacing $U$ by $W$ in \eqref{eq:lec14-dilated-channel} changes the channel by $\mathcal O
(\norm{\mathcal L}_{\mathrm{be}}^2\Delta t^2)$ in diamond norm. Tracing out the ancilla kills all
odd powers of $\sqrt{\Delta t}$, and a direct expansion of $\Tr_a(W(\proj0\otimes\rho)W\dg)$
yields $\rho+\Delta t\,\mathcal L(\rho)+\mathcal O(\norm{\mathcal L}_{\mathrm{be}}^2\Delta t^2)$,
which matches $\ee^{\mathcal L\Delta t}(\rho)=\rho+\Delta t\,\mathcal L(\rho)+\mathcal O
(\norm{\mathcal L}_{\mathrm{be}}^2\Delta t^2)$. Subtracting gives the one-step bound; the
long-time estimate is the telescoping (linear error growth) argument of \cref{thm:lec14-cptp}.
\end{proof}

This construction, illustrated schematically as a chain of \emph{repeated interactions}---couple
to a fresh ancilla, evolve for $\sqrt{\Delta t}$, discard---turns any Hamiltonian simulation
method of \cref{ch:lec11,ch:lec12} into a Lindblad simulator.

\begin{example}[{Ising evolution with dephasing}]
\label{exa:lec14-tfim}
For the transverse-field Ising model $H=-\sum_{j=1}^{n-1}Z_jZ_{j+1}-g\sum_{j=1}^n X_j$ with
site-local dephasing $L_j=\sqrt\gamma\,Z_j$, the Lindbladian is $\dot\rho=-\ii[H,\rho]+\gamma
\sum_j(Z_j\rho Z_j-\rho)$. The dilated Hamiltonian \eqref{eq:lec14-dilated} uses
$\lceil\log_2(n+1)\rceil$ ancilla qubits with the $L_j=\sqrt\gamma Z_j$ in its first
block-column. With $g,\gamma=\mathcal O(1)$ one has $\norm{\mathcal L}_{\mathrm{be}}=\mathcal O(n)$,
so error $\epsilon$ at time $T$ needs $\Delta t=\mathcal O(\epsilon/(T\norm{\mathcal L}
_{\mathrm{be}}^2))$ and $N_t=\mathcal O(T^2 n^2/\epsilon)$ dilated-Hamiltonian simulation steps.
\end{example}

\paragraph{Oblivious amplitude amplification of channels.}
Each dilated step succeeds only when the ancilla postselects back to $\ket0$; naively this
success probability compounds and decays, exactly the pathology of \cref{sec:lec13-3}. The cure
generalizes oblivious amplitude amplification (\cref{ch:lec12}) from unitaries to \emph{channels}.
Given a channel $\mathcal Q(\rho)=\sum_{m}K_m\rho K_m\dg$ with $M=2^a$ Kraus operators and a
\textsc{select} oracle $U_{\textsc{sel}}=\sum_m\proj m\otimes U_{K_m}$, where $U_{K_m}\in
\mathrm{BE}_\alpha(K_m)$, the operator $W=U_{\textsc{sel}}(\Hgate^{\otimes a}\otimes\id)$ is a
block encoding of the Stinespring isometry $V=\sum_m\ket m\otimes K_m$ (\cref{thm:lec02-stinespring}),
with $(\bra{0^b}\otimes\id)W(\ket{0^a}\ket\psi\ket{0^b})=(\alpha\sqrt M)^{-1}V\ket\psi$. Since $V$
is an \emph{isometry}, all its nonzero singular values equal $1$, so amplitude amplification on
the singular value $(\alpha\sqrt M)^{-1}$ is exact: choosing $(\alpha\sqrt M)^{-1}=\sin\frac\pi
2(2k+1)$ maps it to $1$ and eliminates the postselection, yielding a unitary that realizes
$\mathcal Q$ deterministically on the traced-out ancilla. This removes the exponential
success-probability loss from the repeated-interaction chain.

\section{Dissipative state preparation}
\label{sec:lec14-3}

\paragraph{Fixed points as targets.}
A state $\sigma$ is a \emph{fixed point} (stationary state) of the dynamics if $\mathcal L(\sigma)
=0$. For closed systems any $\sigma$ commuting with $H$ is fixed---far too many to be useful. The
dissipative part breaks this degeneracy: a well-chosen $\mathcal L$ can have a \emph{unique} fixed
point $\sigma$, in which case the semigroup is \emph{ergodic} and $\ee^{\mathcal L t}(\rho_0)\to
\sigma$ for every initial $\rho_0$. This is the engine of \emph{dissipative state preparation}:
design $H$ and the $L_j$ so that the target (a thermal state $\ee^{-\beta H_0}/Z$ or a ground
state) is that unique fixed point, then run the dynamics until it mixes. It is the open-system
counterpart of the adiabatic and Zeno solvers of \cref{ch:lec10}: instead of slowly deforming a
Hamiltonian whose ground state is tracked, one relaxes irreversibly toward a designed steady
state.

\paragraph{Building spectrally selective jump operators.}
The construction never diagonalizes $H$. From simple coupling operators $A_a$ (e.g.\ Paulis) one
forms \emph{filtered} jump operators by convolving the Heisenberg-evolved coupling against a
filter $f_\omega$,
\begin{equation}
  L_a(\omega)=\int_{-\infty}^{\infty}f_\omega(s)\,\ee^{\ii Hs}A_a\,\ee^{-\ii Hs}\,\dd s,
  \label{eq:lec14-jump-filter}
\end{equation}
which, in the eigenbasis $H\ket{\psi_i}=\lambda_i\ket{\psi_i}$, acts as $L_a(\omega)=\sum_{i,k}
\hat f_\omega(\lambda_i-\lambda_k)\braket{\psi_i|A_a|\psi_k}\ket{\psi_i}\!\bra{\psi_k}$, with
$\hat f_\omega$ the Fourier transform of $f_\omega$. Thus the filter $\hat f_\omega$ selects which
energy \emph{transitions} the jump operator drives. Because \eqref{eq:lec14-jump-filter} is an
integral of Hamiltonian evolutions, a block encoding of $L_a(\omega)$ is assembled by truncating,
discretizing, and applying LCU (\cref{ch:lec07}) to the Hamiltonian simulations of
\cref{ch:lec11,ch:lec12}.

\begin{proposition}[Ground state as a dark state]
\label{prop:lec14-darkstate}
Let $H$ have a nondegenerate ground state $\ket{\psi_0}$ with gap $\Delta=\lambda_1-\lambda_0>0$,
and take a single coupling $A$ and filter with $\hat f(\nu)=1$ for $\nu\le-\Delta$ and
$\hat f(\nu)=0$ for $\nu\ge0$. Then the jump operator $L=\sum_{i,k}\hat f(\lambda_i-\lambda_k)
\braket{\psi_i|A|\psi_k}\ket{\psi_i}\!\bra{\psi_k}$ only lowers energy, so $L\ket{\psi_0}=0$, and
$\proj{\psi_0}$ is a fixed point of the Lindbladian \eqref{eq:lec14-lindblad} with $H$ as its
coherent term.
\end{proposition}

\begin{proof}
The support condition on $\hat f$ makes every surviving term $\ket{\psi_i}\!\bra{\psi_k}$ carry
$\lambda_i\le\lambda_k-\Delta<\lambda_k$, so $L$ maps each eigenstate strictly downward in energy
and annihilates the lowest, $L\ket{\psi_0}=0$; hence $L\dg L\ket{\psi_0}=0$ and $L\proj{\psi_0}
L\dg=0$. Since also $[H,\proj{\psi_0}]=0$ ($\ket{\psi_0}$ is an eigenstate), every term of
\eqref{eq:lec14-lindblad} annihilates $\proj{\psi_0}$, so $\mathcal L(\proj{\psi_0})=0$. Thus
$\ket{\psi_0}$ is a \emph{dark state}: population flows down the spectrum and pools in the ground
state, which cannot decay further.
\end{proof}

Uniqueness of the dark state (and hence convergence) requires that the coupling connect every
excited state to some lower state---an irreducibility condition on $A$---and quantitative
convergence is measured by the \emph{mixing time}
\begin{equation}
  \tau_{\mathrm{mix}}(\eta)=\inf\Bigl\{t:\ \sup_{\rho_0}\norm{\ee^{\mathcal L t}(\rho_0)-\sigma}_1
    \le\eta\Bigr\}.
  \label{eq:lec14-mixing}
\end{equation}
On $n$ sites one speaks of \emph{fast mixing} when $\tau_{\mathrm{mix}}=\mathcal O(\poly(n))$ and
\emph{rapid mixing} when $\tau_{\mathrm{mix}}=\mathcal O(\polylog(n))$; proving such bounds (the
open-system analogue of the classical Markov-chain spectral-gap estimates) is problem-specific and
an active research area. The total simulation cost is then $\tau_{\mathrm{mix}}$ multiplied by the
per-time cost of one of the simulators of \cref{sec:lec14-2,sec:lec14-4}.

\paragraph{Unraveling.}
For intuition and for classical numerics it helps to represent the same semigroup as an average
over stochastic pure-state trajectories, the \emph{quantum-jump} (Monte Carlo wavefunction)
unraveling: the It\^o stochastic Schr\"odinger equation
$\dd\ket{\psi_t}=\bigl(-\ii H-\tfrac12\sum_j L_j\dg L_j\bigr)\ket{\psi_t}\,\dd t+\sum_j L_j
\ket{\psi_t}\,\dd W_t^j$ satisfies $\E[\proj{\psi_t}]=\rho(t)$ whenever $\E[\proj{\psi_0}]=\rho(0)$.
Between jumps the state evolves under the non-Hermitian effective Hamiltonian $H_{\mathrm{eff}}
=H-\tfrac\ii2\sum_j L_j\dg L_j$; this is exactly the dissipative non-unitary dynamics of
\cref{ch:lec13}, and it motivates the truncated Dyson method below.

\paragraph{Operator splitting.}
An alternative to dilation splits the generator $\mathcal L=\mathcal L_1+\mathcal L_2$ into two
groups (typically coherent versus dissipative) and applies the first-order product
$\Phi(\Delta t)=\ee^{\mathcal L_1\Delta t}\ee^{\mathcal L_2\Delta t}$. The error is governed by
the same superoperator-commutator bound as the Trotter formula (\cref{sec:lec11-1}):
\begin{equation}
  \norm{\ee^{\mathcal L\Delta t}-\ee^{\mathcal L_1\Delta t}\ee^{\mathcal L_2\Delta t}}_\diamond
  \le\tfrac12\Delta t^2\,\norm{[\mathcal L_1,\mathcal L_2]}_\diamond,
  \label{eq:lec14-split}
\end{equation}
so if $[\mathcal L_1,\mathcal L_2]=0$ the splitting is exact. With the coarse bounds
$\norm{\mathcal L_H}_\diamond\le2\norm H$ and $\norm{\mathcal D_j}_\diamond\le2\norm{L_j}^2$ for
the single-jump generator $\mathcal D_j(\rho)=L_j\rho L_j\dg-\tfrac12\{L_j\dg L_j,\rho\}$, the
long-time error is $\mathcal O(t\norm{\mathcal L}_{\mathrm{be}}^2\Delta t)$ as in
\eqref{eq:lec14-dilation-error}. Crucially, unlike the Hamiltonian case, \emph{higher-order}
product formulas cannot be used na\"ively: they require negative time steps, but $\ee^{-\mathcal L
t}$ need not be a channel, so one may not run dissipation backwards. Splitting is attractive when
the pieces are exactly solvable: for site-local dephasing $\mathcal D_j$ commutes with
$\mathcal D_k$ and $\ee^{\mathcal D_j\Delta t}(\rho)=\tfrac{1+\ee^{-2\gamma\Delta t}}2\rho+
\tfrac{1-\ee^{-2\gamma\Delta t}}2 Z_j\rho Z_j$ is an elementary single-qubit dephasing channel.

\section{Nonlinear PDEs via Carleman linearization}
\label{sec:lec14-4}

\paragraph{From nonlinear to infinite-linear.}
Quantum evolution is linear, so a nonlinear differential equation must be recast as a linear one
before a quantum solver applies. Consider the quadratic autonomous ODE---the generic form after
semi-discretizing a quadratic PDE such as Navier--Stokes or a reaction--diffusion equation---
\begin{equation}
  \frac{\dd z}{\dd t}=F_1 z+F_2\,(z\otimes z),\qquad z(0)=z_0\in\C^d,
  \label{eq:lec14-quadratic}
\end{equation}
with $F_1\in\C^{d\times d}$ and $F_2\in\C^{d\times d^2}$. \emph{Carleman linearization} lifts this
to a linear system on the tower of tensor powers $z^{\otimes k}$. Introduce the variables
$y_k:=z^{\otimes k}\in\C^{d^k}$ (so $y_1=z$) and differentiate using the Leibniz rule and
\eqref{eq:lec14-quadratic}:
\begin{equation}
  \frac{\dd y_k}{\dd t}
  =\sum_{\ell=1}^{k}z^{\otimes(\ell-1)}\otimes\dot z\otimes z^{\otimes(k-\ell)}
  =A_k^k\,y_k+A_k^{k+1}\,y_{k+1},
  \label{eq:lec14-carleman-block}
\end{equation}
where $A_k^k$ collects the $k$ copies of the \emph{linear} part $F_1$ (each acting on one tensor
slot, the others carried by identities) and $A_k^{k+1}$ collects the $k$ copies of the
\emph{quadratic} part $F_2$, which raises the level from $k+1$ to $k$. Explicitly,
\begin{equation}
  A_k^k=\sum_{\ell=1}^{k}\id^{\otimes(\ell-1)}\otimes F_1\otimes\id^{\otimes(k-\ell)},
  \qquad
  A_k^{k+1}=\sum_{\ell=1}^{k}\id^{\otimes(\ell-1)}\otimes F_2\otimes\id^{\otimes(k-\ell)}.
  \label{eq:lec14-carleman-coeffs}
\end{equation}
Equation \eqref{eq:lec14-carleman-block} is an \emph{exact} but infinite linear system: the
$k$-th block couples only to itself and the next level $k+1$, an upper-bidiagonal (in block
index) structure. Truncating at level $N$---setting $y_{N+1}\equiv0$---gives a finite linear ODE
\begin{equation}
  \frac{\dd}{\dd t}\hat y=\mathcal A\,\hat y,\qquad
  \hat y=(y_1;y_2;\dots;y_N)\in\C^{d+d^2+\dots+d^N},
  \label{eq:lec14-carleman-trunc}
\end{equation}
whose block-bidiagonal generator $\mathcal A$ carries $A_k^k$ on the diagonal and $A_k^{k+1}$ on
the super-diagonal. This is precisely a linear ODE of the form treated in \cref{ch:lec13}: solve
it with the linear-system, time-marching, LCHS, or Schr\"odingerization method, then read off
$z(T)$ from the first block $y_1$.

\paragraph{Convergence and the dissipativity condition.}
The truncation is useful only if the neglected levels contribute negligibly. This is controlled
by a dimensionless ratio comparing the nonlinear forcing to the linear dissipation. Assume $F_1$
is diagonalizable with eigenvalues of negative real part, order them $\mathrm{Re}\,\lambda_d\le
\dots\le\mathrm{Re}\,\lambda_1<0$, and set
\begin{equation}
  R:=\frac{\norm{z_0}\,\norm{F_2}}{\abs{\mathrm{Re}\,\lambda_1}}.
  \label{eq:lec14-R}
\end{equation}
$R$ is the ratio of the quadratic strength $\norm{F_2}\norm{z_0}$ to the linear decay rate
$\abs{\mathrm{Re}\,\lambda_1}$; it plays the role of a Reynolds number.

\begin{theorem}[Carleman truncation error]
\label{thm:lec14-carleman}
Suppose $F_1$ is diagonalizable with $\mathrm{Re}\,\lambda_1<0$ and $R<1$ as in
\eqref{eq:lec14-R}. Let $z(t)$ solve \eqref{eq:lec14-quadratic} and let $\hat y(t)$ solve the
level-$N$ truncation \eqref{eq:lec14-carleman-trunc} with the exact initial data
$y_k(0)=z_0^{\otimes k}$. Then the error in the recovered first block obeys
\begin{equation}
  \max_{0\le t\le T}\norm{y_1(t)-z(t)}
  \le\norm{z_0}\,R^{N}\,\bigl(1-\ee^{-\abs{\mathrm{Re}\,\lambda_1}t}\bigr)^{N}
  \le\norm{z_0}\,R^{N},
  \label{eq:lec14-carleman-err}
\end{equation}
so the truncation error decays \emph{geometrically} in the level $N$; achieving accuracy
$\epsilon$ needs only $N=\mathcal O\!\bigl(\log(\norm{z_0}/\epsilon)/\log(1/R)\bigr)$.
\end{theorem}

\begin{proof}[Proof sketch]
Diagonalizing $F_1$ and passing to the eigenbasis, the linear part contracts each mode at rate at
least $\abs{\mathrm{Re}\,\lambda_1}$, while the quadratic coupling $A_k^{k+1}$ feeds level $k+1$
into level $k$ with strength $\le k\norm{F_2}$. Writing the error $\eta_k(t)=y_k(t)-z(t)^{\otimes
k}$ (with $\eta_k$ zero above the truncation contributing the sole inhomogeneity at level $N$
through the dropped $A_N^{N+1}y_{N+1}$ term), Duhamel's principle (\cref{prop:lec13-duhamel})
bounds $\norm{\eta_k}$ by a convolution of the decaying propagator against the neglected term.
The neglected term at the top level is $\norm{A_N^{N+1}}\norm{y_{N+1}}\le N\norm{F_2}
\norm{z(t)}^{N+1}$, and $\norm{z(t)}\le\norm{z_0}\ee^{-\abs{\mathrm{Re}\,\lambda_1}t}$ under
$R<1$ (the quadratic term cannot overturn the linear decay). Solving the resulting recursion
$\norm{\eta_1}\lesssim\bigl(\norm{F_2}\norm{z_0}/\abs{\mathrm{Re}\,\lambda_1}\bigr)^{N}\norm{z_0}
=R^N\norm{z_0}$, with the extra $(1-\ee^{-\abs{\mathrm{Re}\,\lambda_1}t})^N$ factor from the
finite-time integrals, gives \eqref{eq:lec14-carleman-err}. The condition $R<1$ is what makes the
geometric series converge; when $R\ge1$ the nonlinear cascade can dominate and the truncation is
not controlled.
\end{proof}

\begin{remark}[Cost and caveats]
\label{rem:lec14-carleman-cost}
The truncated system \eqref{eq:lec14-carleman-trunc} has dimension $\sum_{k=1}^N d^k=\mathcal O
(d^N)$, so it uses $\mathcal O(N\log d)$ qubits---exponentially smaller than the classical
$\mathcal O(d^N)$ state, provided the block-encoding oracles for $\mathcal A$ are efficient (they
inherit the sparsity of $F_1,F_2$). Combined with the linear-ODE solvers of \cref{ch:lec13}, the
end-to-end cost is $\poly(N,d,\alpha_{\mathcal A}T,\log(1/\epsilon))$ and, because $N=\mathcal O
(\log(1/\epsilon)/\log(1/R))$, polylogarithmic in $1/\epsilon$. The essential restriction is
$R<1$: strongly nonlinear (high-Reynolds) regimes fall outside the guarantee. An alternative for
scalar conservation laws and Hamilton--Jacobi equations is the \emph{level-set} approach, which
lifts the nonlinear equation to a \emph{linear} Liouville (transport) equation in phase space at
the cost of extra dimensions; we do not develop it here.
\end{remark}

\section{Example: Burgers' equation via Cole--Hopf}
\label{sec:lec14-5}

\paragraph{An exact linearization.}
Some nonlinear PDEs linearize \emph{exactly}, with no truncation. The archetype is the viscous
Burgers equation, the simplest model combining nonlinear advection with diffusion,
\begin{equation}
  \partial_t u+u\,\partial_x u=\nu\,\partial_{xx}u,\qquad \nu>0,
  \label{eq:lec14-burgers}
\end{equation}
where $\nu$ is the viscosity. The \emph{Cole--Hopf transformation} substitutes
\begin{equation}
  u=-2\nu\,\frac{\partial_x\phi}{\phi}=-2\nu\,\partial_x(\ln\phi),
  \label{eq:lec14-colehopf}
\end{equation}
for a strictly positive potential $\phi(x,t)>0$, and turns \eqref{eq:lec14-burgers} into the
linear heat equation.

\begin{theorem}[Cole--Hopf]
\label{thm:lec14-colehopf}
If $\phi(x,t)>0$ solves the heat equation $\partial_t\phi=\nu\,\partial_{xx}\phi$, then
$u=-2\nu\,\partial_x\phi/\phi$ solves the viscous Burgers equation \eqref{eq:lec14-burgers}.
\end{theorem}

\begin{proof}
Write $u=-2\nu(\ln\phi)_x$. We compute the three terms of \eqref{eq:lec14-burgers} directly. For
the spatial derivatives,
\[
  u_x=-2\nu(\ln\phi)_{xx}=-2\nu\Bigl(\frac{\phi_{xx}}{\phi}-\frac{\phi_x^2}{\phi^2}\Bigr),
  \qquad
  u_{xx}=-2\nu(\ln\phi)_{xxx}
        =-2\nu\Bigl(\frac{\phi_{xxx}}{\phi}-3\frac{\phi_x\phi_{xx}}{\phi^2}
          +2\frac{\phi_x^3}{\phi^3}\Bigr).
\]
The nonlinear term is
\[
  u\,u_x=\Bigl(-2\nu\frac{\phi_x}{\phi}\Bigr)(-2\nu)
    \Bigl(\frac{\phi_{xx}}{\phi}-\frac{\phi_x^2}{\phi^2}\Bigr)
   =4\nu^2\Bigl(\frac{\phi_x\phi_{xx}}{\phi^2}-\frac{\phi_x^3}{\phi^3}\Bigr).
\]
For the time derivative, $u_t=-2\nu(\ln\phi)_{xt}=-2\nu\bigl(\phi_{xt}/\phi-\phi_x\phi_t/\phi^2
\bigr)$. Now use the hypothesis $\phi_t=\nu\phi_{xx}$, hence $\phi_{xt}=\nu\phi_{xxx}$, to get
\[
  u_t=-2\nu^2\Bigl(\frac{\phi_{xxx}}{\phi}-\frac{\phi_x\phi_{xx}}{\phi^2}\Bigr).
\]
Adding the advection term,
\[
  u_t+u\,u_x
  =-2\nu^2\frac{\phi_{xxx}}{\phi}+2\nu^2\frac{\phi_x\phi_{xx}}{\phi^2}
   +4\nu^2\frac{\phi_x\phi_{xx}}{\phi^2}-4\nu^2\frac{\phi_x^3}{\phi^3}
  =-2\nu^2\frac{\phi_{xxx}}{\phi}+6\nu^2\frac{\phi_x\phi_{xx}}{\phi^2}
   -4\nu^2\frac{\phi_x^3}{\phi^3}.
\]
On the other hand,
\[
  \nu\,u_{xx}=-2\nu^2\Bigl(\frac{\phi_{xxx}}{\phi}-3\frac{\phi_x\phi_{xx}}{\phi^2}
    +2\frac{\phi_x^3}{\phi^3}\Bigr)
  =-2\nu^2\frac{\phi_{xxx}}{\phi}+6\nu^2\frac{\phi_x\phi_{xx}}{\phi^2}
   -4\nu^2\frac{\phi_x^3}{\phi^3}.
\]
The two expressions coincide, so $u_t+u u_x=\nu u_{xx}$, as claimed.
\end{proof}

\begin{remark}[Why the substitution is guessed]
\label{rem:lec14-why-colehopf}
The choice \eqref{eq:lec14-colehopf} is natural in two steps. First set $u=\partial_x\psi$
(assuming $u$ has a potential); Burgers becomes $\partial_x(\psi_t+\tfrac12\psi_x^2-\nu\psi_{xx})
=0$, so up to an $x$-independent function $\psi_t+\tfrac12\psi_x^2=\nu\psi_{xx}$, a
Hamilton--Jacobi equation. Second, the exponential (Hopf--Cole) substitution $\psi=-2\nu\ln\phi$
converts the quadratic $\tfrac12\psi_x^2$ into a linear term: $\psi_t=-2\nu\phi_t/\phi$,
$\psi_x=-2\nu\phi_x/\phi$, $\psi_{xx}=-2\nu(\phi_{xx}/\phi-\phi_x^2/\phi^2)$, and substituting
collapses everything to $\phi_t=\nu\phi_{xx}$. The nonlinearity is absorbed because the same
$\phi_x^2/\phi^2$ term appears in $\tfrac12\psi_x^2$ and in $\nu\psi_{xx}$ with matching
coefficients.
\end{remark}

\paragraph{The quantum algorithm.}
Cole--Hopf reduces solving Burgers to three classical/quantum steps. (i) Transform the initial
data: given $u(x,0)=u_0(x)$, set $\phi(x,0)=\exp\bigl(-\tfrac1{2\nu}\int_0^x u_0(x')\,\dd x'\bigr)
>0$, which is exactly the inversion of \eqref{eq:lec14-colehopf} at $t=0$. (ii) Solve the linear
heat equation $\partial_t\phi=\nu\partial_{xx}\phi$ to time $T$ by any method of \cref{ch:lec13}
(linear-system, LCHS, or Schr\"odingerization); after spatial discretization this is the
tridiagonal Laplacian system of \cref{sec:lec13-6} with $\alpha_A=\mathcal O(\nu N^2)$. (iii)
Recover $u(x,T)=-2\nu\,\partial_x\phi/\phi$, a division and a derivative on the grid. Because
$\phi>0$ stays strictly positive (the heat semigroup preserves positivity, and a positive initial
$\phi$ stays positive), the recovery is numerically well posed. The nonlinear solve thus inherits
the near-optimal $\widetilde\bigo(q\,\nu N^2 T)$ scaling of the linear heat-equation solver, with
the norm-growth factor $q$ now referring to the potential $\phi$.

\begin{remark}[Scope]
\label{rem:lec14-scope}
Cole--Hopf is special to Burgers (and a handful of related equations); general quadratic PDEs
have no exact linearization and require the truncated Carleman route of \cref{sec:lec14-4}, whose
guarantee is conditional on $R<1$. Together the two examples bracket the state of the art:
\emph{exact} linearization where a lucky change of variables exists, and \emph{approximate}
linearization with controlled error in the weakly nonlinear regime otherwise.
\end{remark}

\exercisesection

\begin{exercise}\label{ex:lec14-1}
\emph{Trace preservation and the adjoint.} (a) Show directly from \eqref{eq:lec14-lindblad} that
$\Tr[\mathcal L(\rho)]=0$ for every $\rho$, so the Lindblad flow preserves the trace. (b) Verify
that $\mathcal L\dg$ in \eqref{eq:lec14-adjoint} is the Hilbert--Schmidt adjoint of $\mathcal L$,
i.e.\ $\Tr[\mathcal L(X)\dg Y]=\Tr[X\dg\mathcal L\dg(Y)]$ for all $X,Y$, and confirm that
$\mathcal L\dg(\id)=0$ (unitality of the Heisenberg dynamics).
\end{exercise}

\begin{exercise}\label{ex:lec14-2}
\emph{Classical Markov chains.} Complete the computation of \cref{exa:lec14-markov}: with $H=0$
and $L_{ij}=\sqrt{q_{ij}}\,\ket i\!\bra j$, substitute into \eqref{eq:lec14-lindblad} and show
that a diagonal $\rho=\sum_j p_j\proj j$ stays diagonal with $\dot p=Qp$, identifying the entries
of $Q$. Verify that each column of $Q$ sums to zero, so $\ee^{Qt}$ is column-stochastic.
\end{exercise}

\begin{exercise}\label{ex:lec14-3}
\emph{The dilated Hamiltonian reproduces $\mathcal L$.} For a single jump operator $L_1$, take
$\widetilde H$ from \eqref{eq:lec14-dilated} and the channel \eqref{eq:lec14-dilated-channel}.
Expand $U=\ee^{-\ii\widetilde H\sqrt{\Delta t}}$ to second order, trace out the ancilla, and show
$\Tr_a(U(\proj0\otimes\rho)U\dg)=\rho+\Delta t\,\mathcal L(\rho)+\mathcal O(\Delta t^2)$. Explain
why the odd powers of $\sqrt{\Delta t}$ vanish under the partial trace, and why the factor
$H/\sqrt{\Delta t}$ in \eqref{eq:lec14-dilated} is needed for the coherent term to appear at order
$\Delta t$.
\end{exercise}

\begin{exercise}\label{ex:lec14-4}
\emph{Photon loss.} For the cavity model of \cref{exa:lec14-photon} with $H=\omega a\dg a$ and
$L=\sqrt\kappa\,a$, use the adjoint Lindbladian to compute $\mathcal L\dg(n)$ for $n=a\dg a$ and
deduce $\langle n\rangle(t)=\langle n\rangle(0)\ee^{-\kappa t}$. Then compute $\mathcal L\dg(a)$
and find the decay of the coherence $\langle a\rangle(t)$.
\end{exercise}

\begin{exercise}\label{ex:lec14-5}
\emph{Dark state and fixed point.} Prove \cref{prop:lec14-darkstate}: given the energy-lowering
filter, show $L\ket{\psi_0}=0$ and that every term of the Lindbladian annihilates
$\proj{\psi_0}$, so $\mathcal L(\proj{\psi_0})=0$. Explain what additional structural condition on
the coupling $A$ is needed for $\proj{\psi_0}$ to be the \emph{unique} fixed point, and relate the
convergence rate to the mixing time \eqref{eq:lec14-mixing}.
\end{exercise}

\begin{exercise}\label{ex:lec14-6}
\emph{Carleman linearization.} Starting from the quadratic ODE \eqref{eq:lec14-quadratic}, derive
the block recursion \eqref{eq:lec14-carleman-block}--\eqref{eq:lec14-carleman-coeffs} for
$y_k=z^{\otimes k}$ by differentiating $z^{\otimes k}$ with the product rule and substituting
$\dot z=F_1 z+F_2(z\otimes z)$. Explain why the resulting infinite system is block upper-bidiagonal
in the level index $k$, and why truncating at level $N$ (setting $y_{N+1}=0$) leaves a finite
linear ODE solvable by the methods of \cref{ch:lec13}. State the role of the ratio $R$ of
\eqref{eq:lec14-R} in controlling the truncation error.
\end{exercise}

\begin{exercise}\label{ex:lec14-7}
\emph{Cole--Hopf details.} (a) Re-derive \cref{thm:lec14-colehopf} through the potential
substitution of \cref{rem:lec14-why-colehopf}: set $u=\psi_x$, obtain the Hamilton--Jacobi
equation $\psi_t+\tfrac12\psi_x^2=\nu\psi_{xx}$, then substitute $\psi=-2\nu\ln\phi$ and show it
collapses to $\phi_t=\nu\phi_{xx}$. (b) Verify that the initial-data transform $\phi(x,0)=
\exp\bigl(-\tfrac1{2\nu}\int_0^x u_0\bigr)$ inverts \eqref{eq:lec14-colehopf} at $t=0$, and argue
that $\phi>0$ is preserved by the heat flow so the recovery $u=-2\nu\phi_x/\phi$ is well defined
for all $t$.
\end{exercise}

\readingsection
\begin{itemize}[leftmargin=1.4em]
  \item Lin Lin, \emph{live notes}, Ch.~21 (open quantum systems): the Lindblad equation and the
        proof that $\ee^{\mathcal L t}$ is CPTP, the GKSL characterization, dissipative thermal
        and ground-state preparation, the dilated-Hamiltonian (Stinespring) simulation, oblivious
        amplitude amplification of channels, the operator-splitting method, and the truncated
        Dyson method.
  \item H.-P.~Breuer and F.~Petruccione, \emph{The Theory of Open Quantum Systems}, Oxford
        University Press, 2002, for physical background on Lindblad dynamics, quantum-jump
        unravelings, and Markovian master equations.
  \item G.~Lindblad, \emph{On the generators of quantum dynamical semigroups}, Commun.\ Math.\
        Phys.\ \textbf{48}, 119 (1976); and V.~Gorini, A.~Kossakowski, and E.~C.~G.~Sudarshan,
        \emph{Completely positive dynamical semigroups of $N$-level systems}, J.\ Math.\ Phys.\
        \textbf{17}, 821 (1976), for the original GKSL theorem of \cref{prop:lec14-gksl}.
  \item J.-P.~Liu, H.~\O.~Kolden, H.~K.~Krovi, N.~F.~Loureiro, K.~Trivisa, and A.~M.~Childs,
        \emph{Efficient quantum algorithm for dissipative nonlinear differential equations}, Proc.\
        Natl.\ Acad.\ Sci.\ USA \textbf{118}, e2026805118 (2021) (arXiv:2011.03185), for the
        Carleman-linearization algorithm and the $R<1$ convergence condition of
        \cref{sec:lec14-4}.
  \item J.~D.~Cole, \emph{On a quasi-linear parabolic equation occurring in aerodynamics}, Quart.\
        Appl.\ Math.\ \textbf{9}, 225 (1951); and E.~Hopf, \emph{The partial differential equation
        $u_t+uu_x=\mu u_{xx}$}, Comm.\ Pure Appl.\ Math.\ \textbf{3}, 201 (1950), for the exact
        linearization of Burgers' equation (\cref{sec:lec14-5}); see also L.~C.~Evans,
        \emph{Partial Differential Equations}, 2nd ed., AMS, 2010, \S4.4.
  \item S.~Jin and N.~Liu, \emph{Quantum algorithms for nonlinear partial differential equations},
        Bull.\ Sci.\ Math.\ \textbf{194}, 103457 (2024) (arXiv:2212.00681), for the level-set /
        Liouville (Hamilton--Jacobi) approach mentioned in \cref{rem:lec14-carleman-cost}.
\end{itemize}
